\documentclass[11pt,a4paper]{article}
\usepackage[utf8]{inputenc}
\usepackage[T1]{fontenc}
\usepackage{amsmath,amssymb}
\usepackage{hyperref}
\usepackage{listings}
\usepackage{xcolor}
\usepackage{geometry}
\geometry{margin=2.5cm}

\newcommand\YAMLkeystyle{\color{black}\bfseries}
\newcommand\YAMLvaluestyle{\color{blue}\mdseries}

\lstdefinelanguage{yaml}
{
  keywords={true,false,null,y,n},
  keywordstyle=\color{darkgray}\bfseries,
  basicstyle=\YAMLkeystyle,
  sensitive=false,
  comment=[l]{\#},
  morecomment=[s]{/*}{*/},
  commentstyle=\color{purple}\ttfamily,
  stringstyle=\YAMLvaluestyle\ttfamily,
  moredelim=[l][\color{orange}]{\&},
  moredelim=[l][\color{magenta}]{*},
  moredelim=**[il][\color{purple}]{:},
  morestring=[b]',
  morestring=[b]"
}

\lstset{
  basicstyle=\ttfamily\small,
  backgroundcolor=\color{gray!10},
  frame=single,
  breaklines=true,
  keywordstyle=\color{blue},
  commentstyle=\color{green!50!black},
  stringstyle=\color{red!70!black},
}

\title{Bayesian Evidence Trees with Monte Carlo Uncertainty Propagation\\[0.5em]
\large Theory, Likelihood Estimation, Adversarial Analysis, and Implementation Guide}
\author{Ari-Pekka Sihvonen\\
\texttt{ap.sihvonen@gmail.com}}
\date{\copyright\ 2026 — MIT License}

\begin{document}
\maketitle

\begin{abstract}
This document describes the theory, implementation, and practical application of
Bayes Tree, a framework for evaluating complex hypotheses by decomposing them
into hierarchical evidence trees. Top-level branches act as independent evidence
combined at the root via likelihood ratios in log-odds space, while Monte Carlo
simulation propagates uncertainty from likelihood ratio intervals to the
posterior distribution. We present the formal model, provide systematic methods
for estimating likelihood ratios from data, expert judgement, and the
literature, and document the tool's sensitivity analysis and importance ranking
capabilities. We further introduce an \emph{adversarial robustness analysis}
framework that systematically challenges tree assumptions---prior specification,
likelihood ratio calibration, distributional form, and evidence
independence---to identify vulnerabilities and generate structured defenses.
\end{abstract}

\tableofcontents
\newpage

%═══════════════════════════════════════════════════════════════════════════════
\section{Introduction}

The Bayesian Decision Tree tool combines Bayesian inference with Monte Carlo
simulation to propagate uncertainty through a tree of evidence. Each branch
of the tree represents an independent piece of evidence characterized by a
\emph{likelihood ratio} (LR), and the tool computes a posterior probability
distribution for the root hypothesis.

The motivation is to provide a structured, transparent framework for combining
multiple heuristic likelihood ratios into a posterior probability. It is useful
for exploratory ``what-if'' analysis, sensitivity studies, and producing an
interpretable tree of how local evidence shifts beliefs.

%═══════════════════════════════════════════════════════════════════════════════
\section{Theoretical Background}

\subsection{Bayes' Theorem}

For a hypothesis $H$ and evidence $E$:
\begin{equation}
P(H|E) = \frac{P(E|H)\,P(H)}{P(E)}
\end{equation}

The \textbf{likelihood ratio} (LR) form is:
\begin{equation}
\text{LR} = \frac{P(E|H)}{P(E|\neg H)}
\end{equation}

\subsection{Log-Odds Representation}

Converting probability to log-odds:
\begin{equation}
\text{lo}(p) = \ln\frac{p}{1-p}
\end{equation}

And back:
\begin{equation}
p = \frac{1}{1 + e^{-\text{lo}}}
\end{equation}

A Bayesian update in log-odds space is simply additive:
\begin{equation}
\text{lo}(P(H|E)) = \text{lo}(P(H)) + \ln(\text{LR})
\end{equation}

This additive property makes it natural to combine multiple independent pieces
of evidence.

\subsection{Combining Independent Evidence}

Given $n$ independent evidence branches with likelihood ratios
$\text{LR}_1, \text{LR}_2, \ldots, \text{LR}_n$, the combined posterior
in log-odds is:
\begin{equation}
\text{lo}(\text{posterior}) = \text{lo}(\text{prior}) + \sum_{i=1}^{n} \ln(\text{LR}_i)
\end{equation}

Equivalently, in odds form:
\begin{equation}
\text{odds}(H \mid E_{1:n}) = \text{odds}(H) \prod_{i=1}^{n} \text{LR}_i
\end{equation}

The tool computes this at the root level, treating each top-level child as an
independent evidence source.

\subsection{Effective Likelihood Ratio}

The \emph{effective LR} is the single likelihood ratio that would produce the
same posterior from the given prior:
\begin{equation}
\text{LR}_{\text{eff}} = \exp\!\Big(\text{lo}(\text{posterior}) - \text{lo}(\text{prior})\Big)
\end{equation}

\subsection{Monte Carlo Uncertainty Propagation}

When the LR for each branch is uncertain (specified as an interval rather than
a point value), the tool uses Monte Carlo simulation:

\begin{enumerate}
\item For each simulation run, sample an LR from each branch's distribution.
\item Compute the combined posterior using the log-odds sum.
\item Repeat $N$ times (default: 10{,}000) to build a posterior distribution.
\item Report summary statistics: median, mean, 90\% credible interval.
\end{enumerate}

\subsection{LR Sampling Distributions}

Three distributions are available for sampling LR values from the interval
$[\text{lr\_min}, \text{lr\_max}]$:

\begin{description}
\item[log\_uniform] (default) --- Uniform in log-space. Best when uncertainty
  spans orders of magnitude. The sampled value is:
  \begin{equation}
  \text{LR} = \exp\!\Big(\text{Uniform}\big(\ln(\text{lr\_min}),\, \ln(\text{lr\_max})\big)\Big)
  \end{equation}

\item[uniform] --- Uniform in linear space:
  \begin{equation}
  \text{LR} = \text{Uniform}(\text{lr\_min},\, \text{lr\_max})
  \end{equation}

\item[beta] --- Beta distribution scaled to the interval, controlled by
  parameters \texttt{lr\_alpha} and \texttt{lr\_beta}:
  \begin{equation}
  \text{LR} = \text{lr\_min} + \text{Beta}(\alpha, \beta) \times (\text{lr\_max} - \text{lr\_min})
  \end{equation}
\end{description}

%═══════════════════════════════════════════════════════════════════════════════
\section{Tree Structure and Chaining}

The tree has two distinct roles:
\begin{itemize}
\item \textbf{Root level:} Children are combined as independent evidence
  (log-odds sum). Each child contributes its own LR to the root posterior.
\item \textbf{Sub-levels:} Children represent drill-down or decomposition of
  the parent evidence. The child's prior is set to the parent's posterior
  (chaining), showing how the probability evolves through successive reasoning
  steps.
\end{itemize}

Importantly, children of top-level branches are \emph{not} multiplied into the
root combination. They only serve to explain or decompose a branch's LR. If
child-level evidence should aggregate to the root, the model must be
restructured so that the relevant nodes appear at the top level.

%═══════════════════════════════════════════════════════════════════════════════
\section{Sensitivity Analysis}
\label{sec:sensitivity}

The tool provides two sensitivity measures:
\begin{enumerate}
\item \textbf{Threshold probabilities:} $P(\text{posterior} < 5\%)$,
  $P(\text{posterior} < 10\%)$, $P(\text{posterior} > 50\%)$.
\item \textbf{Leave-one-out ranking:} For each branch, compute the posterior
  without that branch. The difference from the baseline median shows how much
  each branch influences the result.
\end{enumerate}

%═══════════════════════════════════════════════════════════════════════════════
\section{YAML File Structure}

The input is a YAML file defining a tree of evidence. Below is the complete
specification.

\subsection{Root Node}

\begin{lstlisting}[language=yaml]
node: "Hypothesis name"
prior: 0.5                  # Prior probability (0 to 1)
children:
  - ...                     # List of evidence branches
\end{lstlisting}

\subsection{Evidence Node}

Each child node represents one piece of evidence:

\begin{lstlisting}[language=yaml]
- node: "Evidence description"
  evidence_type: for        # "for", "against", or "neutral"

  # Option A: Point estimate
  likelihood_ratio: 2.5     # Exact LR value

  # Option B: Uncertainty interval (preferred)
  lr_min: 0.01              # Lower bound of LR
  lr_max: 0.10              # Upper bound of LR
  lr_dist: log_uniform      # Distribution: log_uniform | uniform | beta

  # Optional for beta distribution
  lr_alpha: 2               # Beta shape parameter alpha
  lr_beta: 2                # Beta shape parameter beta

  # Optional sub-evidence (drill-down)
  children:
    - node: "Sub-evidence"
      evidence_type: for
      likelihood_ratio: 1.5
\end{lstlisting}

\subsection{Evidence Type}

\begin{description}
\item[for] --- Evidence supporting the hypothesis ($\text{LR} > 1$).
\item[against] --- Evidence against the hypothesis ($\text{LR} < 1$).
\item[neutral] --- Uninformative evidence ($\text{LR} \approx 1$).
\end{description}

The tool validates that \texttt{evidence\_type} is consistent with the LR
values and warns if there is a mismatch.

\subsection{Likelihood Ratio Interpretation}

\begin{center}
\begin{tabular}{ll}
\hline
\textbf{LR value} & \textbf{Meaning} \\
\hline
LR $\gg$ 1 & Strong evidence for the hypothesis \\
LR $>$ 1   & Moderate evidence for \\
LR $=$ 1   & Neutral (no update) \\
LR $<$ 1   & Evidence against \\
LR $\ll$ 1 & Strong evidence against \\
\hline
\end{tabular}
\end{center}

%═══════════════════════════════════════════════════════════════════════════════
\section{Estimating Likelihood Ratios}
\label{sec:estimating-lr}

The quality of a Bayes Tree analysis depends on the quality of the likelihood
ratio estimates assigned to each branch. This section presents systematic
methods for arriving at defensible LR values, ranging from data-driven
computation to structured expert judgement.

\subsection{From Data: Frequency-Based Estimation}

When empirical data are available, the likelihood ratio can be computed
directly from conditional frequencies:
\begin{equation}
\text{LR} = \frac{P(E \mid H)}{P(E \mid \neg H)}
= \frac{\text{rate of evidence } E \text{ when } H \text{ is true}}
       {\text{rate of evidence } E \text{ when } H \text{ is false}}
\end{equation}

\paragraph{Example: diagnostic test.}
A blood test detects a disease in 95\% of true cases (sensitivity = 0.95) and
produces false positives in 5\% of healthy individuals (1 $-$ specificity = 0.05):
\[
\text{LR}^+ = \frac{0.95}{0.05} = 19.0
\]
A positive test result is strong evidence for disease
($\text{LR} \approx 19$). A negative result gives:
\[
\text{LR}^- = \frac{1 - 0.95}{1 - 0.05} = \frac{0.05}{0.95} \approx 0.053
\]

When exact frequencies are uncertain, use the range of plausible values as
\texttt{lr\_min} and \texttt{lr\_max}. For instance, if sensitivity could be
90--98\% and specificity 93--97\%:
\[
\text{LR}^+_{\min} = \frac{0.90}{0.07} \approx 12.9, \qquad
\text{LR}^+_{\max} = \frac{0.98}{0.03} \approx 32.7
\]

\subsection{From Published Studies}

Many empirical fields report results from which LRs can be derived:

\begin{description}
\item[Relative risk (RR) or odds ratio (OR).] When baseline rates are known,
  $\text{LR} \approx \text{RR}$ for rare outcomes. For case-control studies,
  $\text{OR} \approx \text{LR}$ when the sampling fraction is small.

\item[Effect sizes.] Cohen's $d$ can be converted to an approximate LR for
  observing a value above a threshold:
  \[
  \text{LR} \approx \frac{\Phi(\text{threshold} - d)}{\Phi(\text{threshold})}
  \]
  where $\Phi$ is the standard normal CDF. For a one-tailed observation at
  the median ($\text{threshold} = 0$): $\text{LR} = \Phi(-d/2) / \Phi(d/2)$.

\item[Bayes factors.] If a study reports a Bayes factor $\text{BF}_{10}$, this
  is \emph{already} a likelihood ratio and can be used directly.

\item[Confidence intervals.] A 95\% CI that excludes the null gives a
  rough lower bound on evidential strength.  The further the CI is from the
  null, the stronger the evidence.
\end{description}

\paragraph{Handling uncertainty in published estimates.}
Use the study's confidence interval or reported range to set
\texttt{lr\_min}/\texttt{lr\_max}. If a meta-analysis reports
$\text{OR} = 3.2$ with 95\% CI $[1.8, 5.7]$, assign
$\texttt{lr\_min} = 1.8$, $\texttt{lr\_max} = 5.7$.

\subsection{From Expert Judgement: Structured Elicitation}

When empirical data are sparse or unavailable, likelihood ratios must be
elicited from domain expertise. The following structured methods reduce bias.

\subsubsection{The Surprise Test}

Ask: \emph{``How surprised would I be to see this evidence if $H$ were true,
compared to if $H$ were false?''}

\begin{center}
\begin{tabular}{lrl}
\hline
\textbf{Reaction if $H$ true vs.\ false} & \textbf{Approximate LR} & \textbf{Strength}\\
\hline
Equally likely either way & 1 & Neutral \\
Somewhat more expected if $H$ & 2--3 & Weak \\
Clearly more expected if $H$ & 5--10 & Moderate \\
Very surprising if $\neg H$ & 20--50 & Strong \\
Nearly impossible if $\neg H$ & 100+ & Decisive \\
\hline
Somewhat more expected if $\neg H$ & 0.3--0.5 & Weak against \\
Clearly more expected if $\neg H$ & 0.1--0.2 & Moderate against \\
Very surprising if $H$ & 0.02--0.05 & Strong against \\
\hline
\end{tabular}
\end{center}

\subsubsection{The Frequency Frame}

Convert the question to a natural frequency to reduce cognitive bias:

\emph{``Out of 1000 cases where $H$ is true, how many would show evidence $E$?
Out of 1000 cases where $H$ is false, how many would show $E$?''}

If the answers are 800 and 200:
\[
\text{LR} = \frac{800/1000}{200/1000} = 4.0
\]

\subsubsection{The Reference Class Method}

Anchor the estimate to a comparable situation where data exist:

\begin{enumerate}
\item Identify a structurally similar domain where the same type of evidence
  has been studied empirically.
\item Use that domain's LR as a starting point.
\item Adjust up or down for known differences between the reference and target.
\item Express residual uncertainty as \texttt{lr\_min}/\texttt{lr\_max}.
\end{enumerate}

\paragraph{Example.} You need the LR for ``manuscript attributed to author $X$
by stylometric analysis.'' Stylometric authorship studies in English literature
typically achieve LRs of 5--50 depending on method and corpus. If your case
involves an unusually small corpus: $\texttt{lr\_min} = 2$,
$\texttt{lr\_max} = 15$.

\subsubsection{The Adversarial Test}

For politically or emotionally charged hypotheses, adopt the strongest
plausible version of the opposing view:

\begin{enumerate}
\item Estimate $\text{LR}_{\text{proponent}}$: the LR a reasonable proponent
  of $H$ would assign.
\item Estimate $\text{LR}_{\text{sceptic}}$: the LR a reasonable sceptic
  would assign.
\item Use these as \texttt{lr\_max} and \texttt{lr\_min} respectively.
\end{enumerate}

This ensures the interval spans the range of defensible interpretations.

\subsection{Calibration Rules of Thumb}

\begin{center}
\begin{tabular}{rll}
\hline
\textbf{LR} & \textbf{Verbal label} & \textbf{Typical source} \\
\hline
$> 100$    & Decisive      & Forensic DNA, mathematical proof \\
$20$--$100$ & Strong        & Controlled experiment, multiple witnesses \\
$5$--$20$   & Moderate      & Observational study, expert testimony \\
$2$--$5$    & Weak          & Anecdotal report, weak correlation \\
$1$         & Neutral       & Irrelevant information \\
$0.2$--$0.5$ & Weak against  & Minor disconfirmation \\
$0.05$--$0.2$ & Moderate against & Failed prediction, conflicting data \\
$0.01$--$0.05$ & Strong against & Robust negative result \\
$< 0.01$   & Decisive against & Definitive refutation \\
\hline
\end{tabular}
\end{center}

\subsection{Common Pitfalls}

\begin{enumerate}
\item \textbf{Conflating LR with posterior.} An LR of 10 does \emph{not}
  mean 90\% probability. It shifts the odds by a factor of 10 --- the
  posterior depends on the prior.

\item \textbf{Ignoring base rates.} When converting published statistics to
  LRs, ensure you are computing $P(E|H)/P(E|\neg H)$, not $P(H|E)$.

\item \textbf{Double-counting evidence.} If two branches draw on the same
  underlying data, they are not independent. Merge them or model one as a
  sub-child.

\item \textbf{Overconfidence in point estimates.} Unless the LR is known from
  large-sample data, prefer an interval. Specify \texttt{lr\_min} and
  \texttt{lr\_max} to represent your genuine uncertainty.

\item \textbf{Anchoring on round numbers.} Analysts tend to assign LRs of
  exactly 2, 5, or 10. Consider whether 3.5 or 7 might be more accurate.

\item \textbf{Symmetry bias.} Absence of evidence is \emph{not} evidence of
  absence. If a search for evidence $E$ was conducted and nothing was found,
  the LR depends on the search's power to detect $E$ if $H$ were true.
\end{enumerate}

\subsection{Worked Example: Historical Hypothesis}

Consider the hypothesis ``Napoleon was deliberately poisoned with arsenic.''

\paragraph{Evidence: hair samples contain high arsenic.}
\begin{itemize}
\item $P(E|H)$: If Napoleon was poisoned, high arsenic in hair is very likely.
  Estimate: 0.90--0.99.
\item $P(E|\neg H)$: If not poisoned, Scheele's Green wallpaper, medicines,
  and environmental arsenic could still produce elevated levels. Studies show
  this was common in the 19th century. Estimate: 0.40--0.70.
\item $\text{LR} = P(E|H)/P(E|\neg H)$:
  $\text{lr\_min} = 0.90/0.70 \approx 1.3$,
  $\text{lr\_max} = 0.99/0.40 \approx 2.5$.
\end{itemize}

The evidence is only weakly diagnostic because the alternative explanation
(environmental arsenic) has high $P(E|\neg H)$.

\paragraph{Evidence: autopsy showed stomach cancer.}
\begin{itemize}
\item $P(E|H)$: If poisoned, a stomach ulcer/tumour is possible but not
  expected. Estimate: 0.10--0.30.
\item $P(E|\neg H)$: If not poisoned, stomach cancer explains the death
  directly. Estimate: 0.70--0.90.
\item $\text{LR} = 0.10/0.90 \approx 0.11$ to $0.30/0.70 \approx 0.43$.
\end{itemize}

This is moderate evidence \emph{against} poisoning.

%═══════════════════════════════════════════════════════════════════════════════
\section{Complete YAML Example}

\begin{lstlisting}[language=yaml]
node: "Patient has disease X"
prior: 0.10
children:
  - node: "Positive blood test"
    evidence_type: for
    lr_min: 3.0
    lr_max: 8.0
    lr_dist: log_uniform
    children:
      - node: "Test specificity concerns"
        evidence_type: against
        lr_min: 0.5
        lr_max: 0.9

  - node: "Family history present"
    evidence_type: for
    likelihood_ratio: 2.0

  - node: "Age below typical onset"
    evidence_type: against
    lr_min: 0.3
    lr_max: 0.7
    lr_dist: uniform

  - node: "Symptom pattern matches"
    evidence_type: for
    lr_min: 1.5
    lr_max: 4.0
    lr_dist: beta
    lr_alpha: 2
    lr_beta: 3
\end{lstlisting}

%═══════════════════════════════════════════════════════════════════════════════
\section{Usage}

\subsection{Running the Tool}

\begin{lstlisting}[language=bash]
# Install
pip install bayes-tree

# Terminal UI with histogram and tree view
bayes-tree examples/mycase.yaml

# Machine-readable JSON output
bayes-tree examples/mycase.yaml --format json

# Prior sensitivity sweep
bayes-tree examples/mycase.yaml --prior-sweep

# Custom number of simulations
bayes-tree examples/mycase.yaml -n 50000

# Adversarial robustness audit
bayes-tree examples/mycase.yaml --adversarial

# Adversarial audit with JSON output
bayes-tree examples/mycase.yaml --adversarial --format json
\end{lstlisting}

\subsection{Output Sections}

The tool produces the following output:
\begin{enumerate}
\item \textbf{Validation warnings} --- Flags any inconsistencies between
  \texttt{evidence\_type} and LR values.
\item \textbf{Combined posterior} --- Summary statistics (median, mean, CI)
  and ASCII histogram of the posterior distribution.
\item \textbf{Tree view} --- Hierarchical display showing each branch's
  contribution, with prior $\to$ posterior transitions.
\item \textbf{Sensitivity analysis} --- Threshold probabilities.
\item \textbf{Importance ranking} --- Leave-one-out analysis showing which
  branches have the most influence.
\item \textbf{Adversarial audit} (with \texttt{-{}-adversarial}) --- Robustness
  report listing plausible attacks, severity ratings, and defense suggestions.
\end{enumerate}

\subsection{Tips for Creating YAML Files}

\begin{enumerate}
\item \textbf{Start with the prior.} What is the base rate or initial
  probability before any evidence is considered?
\item \textbf{List independent evidence.} Each top-level child should be
  logically independent of the others.
\item \textbf{Prefer intervals over point estimates.} Using \texttt{lr\_min}
  and \texttt{lr\_max} captures your uncertainty about the strength of
  evidence.
\item \textbf{Use log\_uniform for uncertain LRs.} When you know the LR is
  ``somewhere between 0.01 and 0.1'' but aren't sure where, log-uniform
  gives appropriate weight to different orders of magnitude.
\item \textbf{Add children for drill-down.} Sub-nodes explain or decompose
  a parent's evidence but don't add to the root combination.
\item \textbf{Validate your intuitions.} Run the tool and check that the
  output matches your qualitative expectation. If not, reconsider your LR
  estimates.
\end{enumerate}

%═══════════════════════════════════════════════════════════════════════════════
\section{Adversarial Robustness Analysis}

A Bayesian evidence tree encodes several assumptions: the prior probability,
the likelihood ratio estimates, the distributional forms, and the independence
of top-level branches.  An \emph{adversarial audit} systematically challenges
each of these assumptions with plausible perturbations, measures the impact on
the posterior, and reports which assumptions the conclusion is most sensitive
to.  This mirrors the role of a sceptical peer reviewer who searches for the
weakest points in an argument.

\subsection{Motivation}

Even a well-constructed evidence tree may be fragile: the conclusion may depend
critically on a single uncertain LR estimate, or collapse if two branches are
correlated rather than independent.  Standard sensitivity analysis
(Section~\ref{sec:sensitivity}) varies one parameter at a time; adversarial
analysis goes further by (i)~testing specific \emph{types} of methodological
challenge, (ii)~scoring each attack's real-world \emph{plausibility}, and
(iii)~searching for \emph{minimal attack combinations} that flip the
conclusion.

\subsection{Attack Types}

The framework implements four attack classes, each targeting a different
modelling assumption.

\subsubsection{Prior Bias Attack}

The prior $P(H)$ is often the most contested parameter.  The prior bias attack
tests robustness by sweeping the prior through a range of values
$\{0.01, 0.05, 0.10, \ldots, 0.99\}$ and reports:
\begin{itemize}
\item The posterior at each perturbed prior.
\item The \emph{flip prior} $p^*$: the threshold prior above (or below) which
  the posterior crosses 50\%, found by binary search:
  \[
    p^* = \arg\min_p \bigl| \widetilde{P}(H \mid E;\, p) - 0.50 \bigr|
  \]
  where $\widetilde{P}$ denotes the median posterior from Monte Carlo
  simulation with prior $p$.
\end{itemize}

\subsubsection{LR Calibration Attack}

This attack questions the likelihood ratio estimates by applying two
perturbations:
\begin{enumerate}
\item \textbf{Shrinkage toward neutrality.}  The log-LR is multiplied by a
  shrinkage factor $\gamma \in \{0.25, 0.50, 0.75\}$:
  \[
    \log \text{LR}' = \gamma \cdot \log \text{LR}
  \]
  This pulls extreme LR values (both strong support and strong counter-evidence)
  toward the neutral value $\text{LR} = 1$.
\item \textbf{Uncertainty expansion.}  The log-space interval
  $[\log \text{LR}_{\min},\, \log \text{LR}_{\max}]$ is expanded by a factor
  $\lambda \in \{1.5, 2.0, 3.0\}$ around the centre:
  \[
    \log \text{LR}'_{\min/\max}
      = c \mp \lambda \cdot \tfrac{w}{2},
    \qquad
    c = \tfrac{\log \text{LR}_{\min} + \log \text{LR}_{\max}}{2},
    \quad
    w = \log \text{LR}_{\max} - \log \text{LR}_{\min}
  \]
\end{enumerate}

\subsubsection{Distribution Misspecification Attack}

The choice of sampling distribution (log-uniform, uniform, or beta) affects the
posterior.  This attack tests the sensitivity to that choice by replacing each
branch's distribution with alternatives:
\begin{itemize}
\item Log-uniform $\to$ uniform, beta($\alpha{=}2, \beta{=}5$), beta($\alpha{=}5, \beta{=}2$)
\item Uniform $\to$ log-uniform, beta variants
\end{itemize}
If the posterior is stable across distributional assumptions, the conclusion is
more trustworthy.

\subsubsection{Correlation Attack}

The standard model assumes independence of top-level branches.  The correlation
attack introduces dependence between a pair of branches using a Gaussian copula:
\begin{enumerate}
\item Draw correlated standard normals $(Z_1, Z_2)$ with correlation $\rho$:
  \[
    \begin{pmatrix} Z_1 \\ Z_2 \end{pmatrix}
    \sim \mathcal{N}\!\left(
      \mathbf{0},\,
      \begin{pmatrix} 1 & \rho \\ \rho & 1 \end{pmatrix}
    \right)
  \]
\item Transform to uniform quantiles via $U_i = \Phi(Z_i)$.
\item Use $U_i$ as the quantile input for each branch's LR sampler.
\end{enumerate}
This is tested for $\rho \in \{0.3, 0.5, 0.7, 0.9\}$.  The attack is most
plausible when two branches share a common information source.

\subsection{Plausibility Scoring}

Not every mathematically possible attack is realistic.  Each attack receives a
\emph{plausibility score} on a 0--10 scale, based on:
\begin{itemize}
\item \textbf{Evidence type.}  Qualitative evidence (testimony, motive,
  circumstantial claims) is more attackable than quantitative evidence (DNA,
  isotope analysis, controlled experiments).
\item \textbf{LR extremity.}  Very large LR values ($> 100$) are harder to
  justify and thus more attackable.
\item \textbf{Correlation context.}  Same-source evidence pairs receive higher
  correlation plausibility.
\item \textbf{Prior shift magnitude.}  Small prior shifts are more plausible
  than extreme ones.
\end{itemize}
Only attacks with plausibility $\geq 3.0$ (configurable) are reported.

\subsection{Attack Search}

Given an \emph{attack budget} $k$ (default $k = 3$, modelling a realistic peer
review), the search algorithm finds the minimal combination of attacks that
flips the conclusion below a threshold (default 50\%).

The greedy algorithm proceeds as follows:
\begin{enumerate}
\item Score each plausible attack by $\text{impact} \times (0.5 + 0.5 \cdot
  \text{plausibility}/10)$.
\item Select the highest-scoring attack, apply it to the tree.
\item Re-evaluate remaining attacks on the modified tree.
\item Repeat up to $k$ times or until the conclusion flips.
\end{enumerate}

\subsection{Severity Classification}

Each attack result is classified by the absolute change in posterior median:
\begin{center}
\begin{tabular}{lll}
\textbf{Severity} & $|\Delta \text{posterior}|$ & \textbf{Interpretation} \\
\hline
Critical & $> 20\%$ & Conclusion may not survive scrutiny \\
High     & $> 10\%$ & Significant vulnerability \\
Moderate & $> 5\%$  & Non-trivial sensitivity \\
Low      & $\leq 5\%$ & Minor impact \\
\end{tabular}
\end{center}

\subsection{Defense Generation}

For each successful attack, the system generates structured defenses:
\begin{itemize}
\item \textbf{Prior bias:} Justify the prior with base rates; argue why priors
  above the flip threshold are unreasonable.
\item \textbf{LR calibration:} Cite specific studies supporting the LR estimate;
  compute LR from documented base rates.
\item \textbf{Misspecification:} Justify the chosen distribution; show
  robustness by running with all three distributions.
\item \textbf{Correlation:} Add corroborating independent evidence; document
  source independence explicitly.
\end{itemize}

\subsection{Interpreting the Audit}

The audit produces a verdict:
\begin{itemize}
\item \textbf{ROBUST:} No high-severity attacks found; the conclusion
  withstands plausible challenges.
\item \textbf{CAUTIOUS:} High-severity attacks exist but do not flip the
  conclusion; additional evidence or justification is recommended.
\item \textbf{VULNERABLE:} The conclusion can be flipped by a plausible
  combination of attacks; the argument requires strengthening.
\end{itemize}

\label{sec:adversarial}

%═══════════════════════════════════════════════════════════════════════════════
\section{Limitations and Assumptions}

\begin{itemize}
\item \textbf{Independence assumption.} Top-level branches are assumed
  conditionally independent given $H$ and $\neg H$. This is a strong
  assumption and should be scrutinized when modeling correlated evidence.
  In practice, if two branches share a common information source, they
  should be merged or one should be made a sub-node of the other.
  The adversarial correlation attack (Section~\ref{sec:adversarial}) can
  quantify the impact of violating this assumption.

\item \textbf{Children do not aggregate to root.} Children are used only for
  chained drill-down display. If child-level evidence should contribute to
  the root posterior, the model must be restructured so that the relevant
  nodes appear at the top level.

\item \textbf{Log-uniform bias.} The default log-uniform sampling biases
  toward multiplicative scales. Users should choose a distribution that
  reflects their actual beliefs about where in the interval the true LR lies.

\item \textbf{Subjective LR estimates.} The tool is only as good as the
  likelihood ratios you assign. It makes subjective reasoning transparent
  and quantitative, but does not eliminate subjectivity.

\item \textbf{Numerical stability.} The implementation clamps probabilities
  and log-odds to safe numerical ranges to avoid overflow, which may
  slightly affect extreme posteriors (very close to 0 or 1).
\end{itemize}

%═══════════════════════════════════════════════════════════════════════════════
\section{Implementation Notes}

\begin{itemize}
\item \textbf{Log-odds stability:} Probabilities are clamped to
  $[10^{-12},\, 1-10^{-12}]$ before log-odds conversion; log-odds are
  clamped to $[-700, 700]$ before exponentiation.
\item \textbf{Sampling:} The same random draw strategy is applied per
  simulation; users can set simulation counts via the command line.
\item \textbf{Tree view:} A recursive collection routine computes node-level
  posteriors by feeding a parent's posterior median as prior into children,
  producing local medians and credible intervals.
\item \textbf{Effective LR for root:} The root node's LR interval is derived
  from the 5th--95th percentiles of the effective LR distribution across all
  Monte Carlo runs.
\end{itemize}

%═══════════════════════════════════════════════════════════════════════════════
\section{Dependencies}

The core engine requires:
\begin{itemize}
\item Python 3.9+
\item \texttt{PyYAML} (\texttt{pip install pyyaml})
\end{itemize}

Optional dependencies for visualisation: \texttt{matplotlib}, \texttt{numpy}.
Optional for PDF reports: \texttt{reportlab}.

\end{document}
